\documentclass[a4paper,8pt]{article}
\usepackage{longtable}
\usepackage{booktabs}
\usepackage{geometry}
\usepackage{array}
\usepackage{calc}
\usepackage{tabularx}
\usepackage{enumitem}
\usepackage{seqsplit}
% Fallbacks for the generated fragments (the paper's LuaLaTeX pipeline
% provides the real colored versions via bayeslfmt.lua)
\providecommand{\bprint}[1]{\texttt{#1}}
\providecommand{\prop}[1]{\texttt{#1}}
\DeclareUnicodeCharacter{00AC}{\ensuremath{\lnot}}
\DeclareUnicodeCharacter{21A6}{\ensuremath{\mapsto}}
\DeclareUnicodeCharacter{2212}{\ensuremath{-}}
\DeclareUnicodeCharacter{2227}{\ensuremath{\land}}
\DeclareUnicodeCharacter{2228}{\ensuremath{\lor}}
\DeclareUnicodeCharacter{2264}{\ensuremath{\leq}}
\DeclareUnicodeCharacter{2265}{\ensuremath{\geq}}
\usepackage[round]{natbib}
\bibliographystyle{plainnat}
\newcolumntype{Y}{>{\raggedright\arraybackslash}X}
\geometry{margin=0.25in}
\renewcommand{\arraystretch}{1.0}
\begin{document}

\begin{center}
\Large \textbf{Bayesian Network Model Checking Results}
\end{center}

\vspace{0.3cm}

\noindent
\textbf{\large win95 (76 variables)}
Windows 95 printer troubleshooting diagnostic model for identifying printing problems \cite{heckerman1995}. The network models causal relationships between printer components, system settings, and observed printing failures. Variables include application status, printer drivers, network configuration, paper availability, and other components. Following, we summarize the meaning of properties and analysis results in textual and tabular form:

\begingroup\raggedright
\begin{itemize}[nosep,leftmargin=*]
\item \prop{infl}: \emph{Structural reasoning.} \bprint{INFL(PrtDriver, Problem1|PTROFFLINE)∧ P(Problem1=No\_Output|PrtDriver=No) ≥0.5} Yes, the printer driver influences the symptom and the conditional failure probability is at least 0.5.
\item \prop{map\_conf}: \emph{MAP and MPE.} \bprint{P(MAP(NetPrint, NetOK|Problem1=No\_Output)|Problem1=No\_Output) >=0.5-0.01} The confidence in the most probable explanation is at least 0.49 (within one percentage point of half).
\item \prop{map\_update}: \emph{MAP and MPE.} \bprint{P(MAP(NetPrint, NetOK|Problem1=No\_Output) [AppOK=Incorrect\_Corrupt↦0.5]↔MAP(NetPrint, NetOK|Problem1=No\_Output))=1} The diagnosis is unchanged with certainty: a local printer with an operational network.
\item \prop{prob\_delta}: \emph{What-if analysis.} \bprint{P(Problem1=No\_Output) [PrtCbl=Loose↦0.5]− P(Problem1=No\_Output)} The failure probability increases by 0.1680.
\item \prop{prob\_offline}: \emph{What-if analysis.} \bprint{P(Problem1=No\_Output|PTROFFLINE=Offline∨ ¬(PrtOn=Yes)) [PrtCbl=Loose↦0.5] [PrtMem=Less\_than\_2Mb↦0.4] ≥0.6} Yes, the probability is at least 0.6 under both parameter updates.
\end{itemize}
\endgroup

\begingroup\footnotesize
\noindent\begin{tabularx}{\columnwidth}{@{}l Y c c@{}}
\hline
\textbf{Property} & \textbf{Result (var-elim)} & \multicolumn{2}{c}{\textbf{Time (s)}} \\
\cline{3-4}
 & & \textbf{var-elim} & \textbf{sampling} \\
\hline
\prop{infl} & \bprint{true} & \textless{}1 & 3 \\
\prop{map\_conf} & \bprint{true} & \textless{}1 & 3 \\
\prop{map\_update} & \bprint{true} & \textless{}1 & \textless{}1 \\
\prop{prob\_delta} & \bprint{0.1680} & \textless{}1 & 6 \\
\prop{prob\_offline} & \bprint{true} & \textless{}1 & 3 \\
\hline
\end{tabularx}
\endgroup\par


\vspace{0.5cm}

\noindent
\textbf{\large pathfinder (109 variables)}
Pathology diagnostic network for lymph-node diseases with 109 variables representing disease classes and histological findings \cite{heckerman1992pathfinder}. The network assists surgical pathologists by modelling relationships between disease hypotheses and cellular features. As done for the previous case studies, we summarize the meaning of properties and analysis results in textual and tabular form:

\begingroup\raggedright
\begin{itemize}[nosep,leftmargin=*]
\item \prop{idp}: \emph{Structural reasoning.} \bprint{P(Fault=Tuberculosis|F1=Present) >0.9∧ IDP(F1, F2|Fault)} Tuberculosis is likely and the findings are conditionally independent.
\item \prop{map}: \emph{MAP and MPE.} \bprint{P(MAP(F1|Fault=Tuberculosis)|Fault=Tuberculosis) >=0.9-0.01} The most probable F1 state covers at least 0.89 of the mass.
\item \prop{prob\_complex}: \emph{Marginal and conditional inference.} \bprint{P(Fault=Tuberculosis|F1=Present∧ F2=Yes) >0.9} With both findings present, tuberculosis is the certain diagnosis.
\end{itemize}
\endgroup

\begingroup\footnotesize
\noindent\begin{tabularx}{\columnwidth}{@{}l Y c c@{}}
\hline
\textbf{Property} & \textbf{Result (var-elim)} & \multicolumn{2}{c}{\textbf{Time (s)}} \\
\cline{3-4}
 & & \textbf{var-elim} & \textbf{sampling} \\
\hline
\prop{idp} & \bprint{true} & \textless{}1 & \textless{}1 \\
\prop{map} & \bprint{true} & \textless{}1 & \textless{}1 \\
\prop{prob\_complex} & \bprint{true} & \textless{}1 & \textless{}1 \\
\hline
\end{tabularx}
\endgroup\par


\vspace{0.5cm}

\noindent
\textbf{\large munin (1041 variables)}
Medical diagnostic network for neuromuscular diseases with 1041 variables representing clinical findings \cite{andreassen1989}. The network models complex probabilistic relationships between symptoms, test results, and disease states. It assists physicians in diagnosing rare neuromuscular conditions. As done for win95, we summarize the meaning of properties and analysis results in textual and tabular form:

\begingroup\raggedright
\begin{itemize}[nosep,leftmargin=*]
\item \prop{combined\_idp}: \emph{Combined verification.} \bprint{P(MAP(DIFFN\_DISTR, DIFFN\_TYPE|R\_LNLW\_MED\_SEV=SEV)|R\_LNLW\_MED\_SEV=SEV) >=0.8∧ IDP(DIFFN\_TIME, DIFFN\_SEV|DIFFN\_PATHO)} The most probable distribution and type of neuropathy cover at least 0.8 of the mass and disease duration is independent of severity given the pathology.
\item \prop{map}: \emph{MAP and MPE.} \bprint{MAP(DIFFN\_PATHO, DIFFN\_TYPE|R\_LNLW\_MED\_SEV=SEV)} The result is specific pathology types.
\item \prop{prob\_cond}: \emph{Combined verification.} \bprint{P(DIFFN\_SEV=MOD|R\_LNLW\_MED\_SEV=SEV) <=0.05∧ INFL(R\_LNLW\_MED\_SEV, R\_LNLW\_APB\_MALOSS)} Moderate severity stays below 5\% and median severity drives APB muscle loss.
\end{itemize}
\endgroup

\begingroup\footnotesize
\noindent\begin{tabularx}{\columnwidth}{@{}l Y c c@{}}
\hline
\textbf{Property} & \textbf{Result (var-elim)} & \multicolumn{2}{c}{\textbf{Time (s)}} \\
\cline{3-4}
 & & \textbf{var-elim} & \textbf{sampling} \\
\hline
\prop{combined\_idp} & \bprint{true} & \textless{}1 & 195 \\
\prop{map} & \bprint{DIFFN\_PATHO=AXONAL, DIFFN\_TYPE=MIXED} & \textless{}1 & 338 \\
\prop{prob\_cond} & \bprint{true} & \textless{}1 & \textless{}1 \\
\hline
\end{tabularx}
\endgroup\par


\vspace{0.5cm}
\bibliography{references,main}
\end{document}
